%==============================================================================
% Appendix A: Notation and Conventions
%==============================================================================
\section{Notation and Conventions}
\label{app:A}

This appendix collects the notation conventions, geometric parameters of the
three topologies, definition of the EC connection, and the computational
pipeline used in this paper.

\subsection{Basic Notation and Conventions}
\label{app:A:notation}

\paragraph{Index conventions.}
\begin{itemize}
  \item 4D spacetime indices: $\mu,\nu = 0,1,2,3$ (0 = internal $\Sone$,
        $1,2,3$ = internal 3D space).
  \item Spatial tensor indices: $i,j = 1,2,3$ (internal 3D space).
  \item Spin-2 quintet label: $A=1,\ldots,5$ (traceless symmetric).
  \item Spin-1 triplet label: $k=0,1,2$ (vector).
  \item Topology index: $\Sthree$, $\Tthree$, $\Nilthree$.
\end{itemize}

\paragraph{Geometric quantities.}
\begin{itemize}
  \item $r$: scale parameter of the internal space (common to $\Sthree$ and
        $\Tthree$).
  \item $R$: scale parameter of $\Nilthree$.
  \item $L$: radius of $\Sone$.
  \item $\kappa$: gravitational constant ($\kappa^2 = 8\pi G$).
\end{itemize}

\paragraph{Volume factors.}
\begin{equation}
  \Vol(\Sthree\times\Sone) = 2\pi^2 Lr^3,
  \quad
  \Vol(\Tthree\times\Sone) = (2\pi)^4 Lr^3,
  \quad
  \Vol(\Nilthree\times\Sone) = (2\pi)^4 LR^3.
\end{equation}

\paragraph{Effective potential.}
After minisuperspace reduction, the effective potential is defined as
\begin{equation}
  \Veff = -\!\left(\frac{R_\EC}{2\kappa^2}
          + \theta\,\mathcal{L}_\NY
          + \alpha\,C^2_\EC\right)\cdot\Vol
\end{equation}
(sign convention: $\Veff > 0$ in the stabilisation direction).

\paragraph{Pontryagin density.}
\begin{equation}
  P = \innerprod{R}{\star R}
    = \frac{1}{2}\varepsilon^{abcd}R_{abef}R_{cd}{}^{ef}.
\end{equation}

\paragraph{Auxiliary MIXING cubic coefficient.}
\begin{equation}
  C_\delta :=
  \left.\frac{\partial^3\Veff}{\partial\delta_0\,\partial\delta_1\,\partial\delta_2}
  \right|_{\delta=0}.
\end{equation}

\subsection{Structure Constants and Connections for the Three Topologies}
\label{app:A:struct}

\subsubsection{$\Sthree\times\Sone$}

Structure constants of $\Sthree$ (in the orthonormal basis $e^i = r\sigma^i/2$
from the left-invariant 1-forms $\sigma^i$):
\begin{equation}
  C^i_{jk}(\Sthree) = \frac{4}{r}\varepsilon_{ijk}
  \qquad (i,j,k=1,2,3).
\end{equation}
\textbf{Important:} The coefficient is $4/r$ (not $2/r$), arising from the
normalisation $e^i = r\sigma^i/2$ (see Appendix~\ref{app:B}).

Connection coefficients by the Koszul formula:
\begin{equation}
  \omega_{abc} = \frac{1}{2}\bigl(C_{abc} + C_{cba} - C_{bac}\bigr).
\end{equation}

Volume-preserving squash+shear ansatz ($\det=1$):
\begin{equation}
  e^0 = r(1+\varepsilon)^{2/3}(1+s)^2\,\sigma^0/2,
  \quad
  e^1 = r(1+\varepsilon)^{2/3}(1+s)^{-2}\,\sigma^1/2,
  \quad
  e^2 = r(1+\varepsilon)^{-4/3}\,\sigma^2/2.
\end{equation}

\subsubsection{$\Tthree\times\Sone$}

Structure constants of $\Tthree = \mathbb{R}^3/\mathbb{Z}^3$:
\begin{equation}
  C^a_{bc}(\Tthree) = 0 \quad \forall\,a,b,c.
\end{equation}

Non-volume-preserving squash+shear ansatz:
\begin{equation}
  e^1 = r(1+\varepsilon)\,\dd x^1,
  \quad
  e^2 = r(1+s)\,\dd x^2,
  \quad
  e^3 = r\,\dd x^3.
\end{equation}
Volume factor: $\Vol(\varepsilon,s) = (2\pi)^4 Lr^3(1+\varepsilon)(1+s)$
($\varepsilon,s$-dependent).

\subsubsection{$\Nilthree\times\Sone$}

Structure constants of $\Nilthree$ (3D Heisenberg manifold):
\begin{equation}
  C^2_{01}(\Nilthree) = \frac{1}{R},
  \quad \text{(all others zero)}.
\end{equation}
We adopt the convention $[E_0,E_1]=(1/R)E_2$ in the body of the paper.
The implementation uses Maurer--Cartan forms $\dd\sigma^2 = -\sigma^0\wedge\sigma^1$
as the basis, so the frame coefficient is
$C[2,0,1] = -(1+\varepsilon)^{-4/3}(1+s)^{-2}/R$.  Both describe the same
Heisenberg structure and differ only by a convention.

Volume-preserving squash+shear ansatz:
\begin{equation}
  e^0 = R\,\sigma^0,
  \quad
  e^1 = R(1+\varepsilon)^{2/3}(1+s)\,\sigma^1,
  \quad
  e^2 = R(1+\varepsilon)^{-2/3}(1+s)^{-1}\,\sigma^2.
\end{equation}
Volume-preservation check:
\begin{equation}
  \det(e^0\wedge e^1\wedge e^2)
  \propto (1+\varepsilon)^{2/3-2/3}(1+s)^{1-1} = 1. \checkmark
\end{equation}

LC background Weyl scalar ($\eta=V=0$):
$C^2_\LC(\Nilthree) = 4/(3R^4) \neq 0$ (non-conformal flatness).

\subsection{EC Connection and Contortion}
\label{app:A:ec}

The Einstein--Cartan connection is the sum of the LC connection
$\Gamma^a{}_{bc}$ and the contortion $K^a{}_{bc}$:
\begin{equation}
  \omega_\EC^a{}_{bc} = \Gamma^a{}_{bc} + K^a{}_{bc}.
\end{equation}
Definition of contortion (from torsion $T^a{}_{bc} = 2\omega^a{}_{[bc]}$):
\begin{equation}
  K_{abc} = \frac{1}{2}\bigl(T_{abc} + T_{bca} - T_{cab}\bigr),
  \qquad
  K_{abc} = -K_{acb}.
\end{equation}

\paragraph{Torsion mode definitions (homogeneous ansatz).}
\begin{itemize}
  \item \textbf{AX (axial scalar)} $\eta$: traceless antisymmetric part of
        torsion.  For $\Sthree$: $T^i{}_{jk}\propto\eta\,\varepsilon_{ijk}$.
  \item \textbf{VT (vector-trace)} $V$: trace part of torsion.  For $\Sthree$:
        $T^i{}_{0j}\propto V\,\delta^i_j$.
  \item \textbf{MX mode}: coexistence of AX and VT ($\eta\neq 0$ and $V\neq 0$).
\end{itemize}

EC Weyl scalar correction in the MX mode (all topologies):
\begin{equation}
  C^2_\EC - C^2_\LC
  = \frac{16V^2\eta^2}{3r^2} \quad (\Sthree, \Tthree),
  \qquad
  = \frac{16V^2\eta^2}{3R^2} \quad (\Nilthree).
\end{equation}

\subsection{Pontryagin Protection Mechanisms}
\label{app:A:pont}

\textbf{Pl\"ucker protection (AX mode):}
When $V=0$, the matrix representation of the Riemann tensor satisfies the
Pl\"ucker relation, and $P\equiv 0$.

\textbf{Pfaffian protection (VT mode):}
When $\eta=0$, $P\equiv 0$ by a Pfaffian-type identity.

\subsection{Computational Pipeline}
\label{app:A:pipeline}

The computations in this paper are performed according to the following
pipeline.

\begin{enumerate}
  \item \textbf{Geometric setup}: Set vielbein $e^i$, structure constants
        $C^a_{bc}$, and volume factors for the three topologies
        (\texttt{script/dppu/topology/}).
  \item \textbf{EC connection}: Construct contortion $K$ from the torsion
        ansatz and compute EC connection $\omega_\EC$
        (\texttt{script/dppu/connection/}).
  \item \textbf{Curvature tensors}: Compute from $\omega_\EC$ the curvature
        tensor, Ricci scalar $R_\EC$, Weyl scalar $C^2_\EC$, and Pontryagin
        density $P$ by SymPy (\texttt{script/dppu/curvature/}).
  \item \textbf{Effective potential}: Expand $\Veff(r,\eta,V;\alpha,\theta)$
        explicitly and obtain the Hessian matrix
        (\texttt{script/dppu/action/}).
  \item \textbf{Mode dictionary}: Check spin-2 quintet (squash, off-diagonal),
        spin-1 triplet, and spin-0 mass spectra via generalised eigenvalue
        problems and stationary-point searches
        (\texttt{script/scripts/paper03ec/}, \texttt{script/scripts/proofs/}).
\end{enumerate}

\paragraph{Script correspondence table.}

\begin{table}[H]
\centering
\small
\begin{tabular}{ll}
\toprule
Content & Script \\
\midrule
$C^2_\EC$ symbolic computation (AX/VT/MX dropout) & \texttt{proofs/ax\_vt\_dropout\_proof.py} \\
$\Tthree$ mode dictionary (squash/shear null) & \texttt{paper03ec/t3\_mode\_dictionary.py} \\
$\Nilthree$ EC mode dictionary & \texttt{paper03ec/nil3\_mode\_dictionary.py} \\
$\Tthree$ MX vacuum search & \texttt{paper03ec/t3\_mx\_vacuum\_search.py} \\
$\Sthree$ MX stationary point search & \texttt{paper03ec/s3\_mx\_vacuum\_search.py} \\
$\Nilthree$ MX stationary point search & \texttt{paper03ec/nil3\_mx\_vacuum\_search.py} \\
$\Nilthree$ EFT at EC false vacuum & \texttt{paper03ec/nil3\_false\_vacuum\_eft.py} \\
$\Sthree$ VT spin-2/1 masses & \texttt{paper03ec/s3\_vt\_spin\_masses.py} \\
$\Nilthree$ quintet splitting & \texttt{paper03ec/nil3\_spin2\_quintet\_splitting.py} \\
$\varepsilon$-$s$ cross-term, 3-topology comparison & \texttt{paper03ec/squash\_shear\_cross\_term.py} \\
Palatini protection ($G_{\eta\eta}=G_{VV}=0$) & \texttt{proofs/palatini\_protection\_proof.py} \\
Auxiliary $C_\delta$ null theorem & \texttt{proofs/c\_delta\_null\_proof.py} \\
Cubic channel structure & \texttt{paper03ec/ec\_cubic\_vertex.py} \\
$\gamma=1/2$ analytic proof & \texttt{proofs/gamma\_scaling\_proof.py} \\
$\Tthree$ flatness null test & \texttt{proofs/t3\_flatness\_null\_test.py} \\
\bottomrule
\end{tabular}
\caption{Correspondence between paper content and scripts.}
\label{tab:scripts}
\end{table}

\subsection{Sanity Checks}
\label{app:A:sanity}

The main propositions of this paper were adopted after passing the following
sanity checks.

\begin{enumerate}
  \item AX/VT dropout: $\Delta C^2=0$ (6 cases, SymPy exact zero).
  \item Palatini protection: $G_{\eta\eta}=G_{VV}=0$ (velocity method, 3
        topologies).
  \item Auxiliary $\delta$-sector: $C_\delta=0$ and $\partial_\alpha C_\delta=0$
        (3 topologies $\times$ 3 backgrounds, SymPy exact).
  \item Palatini universality: VT$-$AX spin-2/1 mass $= 0$ (SymPy exact).
  \item $\gamma=1/2$: $r_0^{\text{analytic}}$ vs $r_0^{\text{numeric}}$ error
        $< 0.001$\% (all 6 points).
  \item Quintet zero mode: $m^2(q_3)=0$ (SymPy; cross-terms also $=0$).
  \item Volume preservation: $S^3$, $\Nilthree$ squash+shear $\det=1$ (SymPy).
  \item Cross-term origin check: $\alpha$-dependence verification of
        $\partial^2 V/\partial\varepsilon\partial s$ ($\Tthree$: independent;
        $\Nilthree$: $\alpha$-dependent).
  \item $\Sthree$ MX stationary-point search: no Hessian-positive-definite full
        stationary vacuum in the dictionary range.
  \item $\Nilthree$ MX stationary-point search: no real stationary point with
        $V\eta\neq 0$ in the $\alpha<0$ branch; only local minimum remaining is
        the $\eta=V=0$ false vacuum.
\end{enumerate}
